\begin{figure}[htbp]
	\centering
	\begin{tikzpicture}[>=Stealth, line width=0.8pt, scale=0.8]
		\draw[thick] (-3.5, 0) ellipse [x radius=2.2, y radius=2.8];
		\node at (-3.5, 3.2) {$\mathbb{R}^n$};
		\draw[thick] (-5.7, 0)
		.. controls (-4.5, 0.5) and (-2.5, -0.5) ..
		(-1.3, 0);
		\node at (-3.5,  2.0) {$\Nc(\bm{A})$};
		\fill (-3.5,  1.2) circle (1.5pt) node[right=2pt] {$\bm{x}_n$};
		\node[gray!80] at (-3.5,  0.5) {$\bm{c}_n = \bm{0}$};
		\fill (-3.5, -0.5) circle (1.5pt) node[right=2pt] {$\bm{x}_c$};
		\fill (-3.5, -1.2) circle (1.5pt) node[right=2pt] {$\bm{c}_c$};
		\node at (-3.5, -2.0) {$\Ic(\bm{A}^\top)$};
		\draw[thick] (3.5, 0) ellipse [x radius=2.2, y radius=2.8];
		\node at (3.5, 3.2) {$\mathbb{R}^m$};
		\draw[thick] (1.3, 0)
		.. controls (2.5, 0.5) and (4.5, -0.5) ..
		(5.7, 0);
		\node at (3.5,  2.0) {$\Ic(\bm{A})$};
		\fill (3.5,  1.2) circle (1.5pt) node[right=2pt] {$\bm{b}_c$};
		\fill (3.5,  0.5) circle (1.5pt) node[right=2pt] {$\bm{y}_c$};
		\node[gray!80] at (3.5, -0.5) {$\bm{b}_n = \bm{0}$};
		\fill (3.5, -1.2) circle (1.5pt) node[right=2pt] {$\bm{y}_n$};
		\node at (3.5, -2.0) {$\Nc(\bm{A}^\top)$};
		\draw[->, thick] (-1.0, 1.5) to[bend left=20]
		node[midway, above=4pt] {$\bm{A}\bm{x}_c = \bm{b}_c$}
		(1.0, 1.5);
		\draw[->, thick] (1.0, -1.5) to[bend left=20]
		node[midway, below=4pt] {$\bm{A}^\top\bm{y}_c = \bm{c}_c$}
		(-1.0, -1.5);
	\end{tikzpicture}
	\caption{Struttura dei problemi primale $\bm{A}\bm{x} = \bm{b}$ e duale
		$\bm{A}^\top\bm{y} = \bm{c}$. Il pedice $c$ denota la componente
		\emph{complementare al nucleo}: $\bm{x}_c \in \Ic(\bm{A}^\top)$,
		$\bm{b}_c \in \Ic(\bm{A})$, $\bm{y}_c \in \Ic(\bm{A})$,
		$\bm{c}_c \in \Ic(\bm{A}^\top)$. Il pedice $n$ denota la componente
		nel nucleo: $\bm{x}_n \in \Nc(\bm{A})$, $\bm{y}_n \in \Nc(\bm{A}^\top)$.
		Il primale ammette soluzione se e solo se $\bm{b}_n = \bm{0}$;
		il duale se e solo se $\bm{c}_n = \bm{0}$.}
	\label{fig:problemi_duali}
\end{figure}
